% GENERATED FILE. Edit canonical lessons, then run npm run generate:books.
% canonical-source-hash: fnv1a64:0c05cb31f076f04e
% canonical-lessons: JA-C191-shimaru, JA-C191-ichiba, JA-C191-sakaya, JA-C191-yakkyoku, JA-C191-mitsukeru

\chapter{Closing Time and Local Shops}
\label{ch:ja-closing-time-and-local-shops}
\hlchaptermodality{191}

\begin{chapteropening}
\textbf{By the end of this chapter:} \emph{I can read closing times and find local shops.}

Complete the last lesson of chapter 191: I can read closing times and find local shops.
\end{chapteropening}

\section[\texorpdfstring{\ja{しまる} (shimaru)}{しまる (shimaru)}]{\ja{しまる} (shimaru) — to close; to be shut}

\label{lesson:JA-C191-shimaru}



\begin{quote}
Before the new one: say the Japanese for to enter, then the Japanese for to be open.
\end{quote}

\subsection*{You'll want to know: \ja{しまる}}
\textbf{\ja{しまる}} — \emph{shimaru} — \textquotedblleft{}to close; to be shut\textquotedblright{}.

\begin{quote}
\textbf{\ja{ゆうびんきょくは} \ja{ごじに} \ja{しまります。}} — \emph{yūbinkyoku wa goji ni shimarimasu} — \textquotedblleft{}The post office closes at five.\textquotedblright{}
\end{quote}

Polite \textbf{\ja{しまります}}, \emph{shimarimasu}; a godan verb.

\begin{cousinweb}
\textbf{\ja{しまる}} (\emph{shimaru}) is a native Japanese word (wago), inherited rather than borrowed from Chinese; its kanji was attached later as a kun'yomi (JMdict, KANJIDIC2).
\end{cousinweb}

\begin{grammarlens}[title={What you've built}]
One more word for everyday life at A2.
\end{grammarlens}

\subsection*{Your turn}
\begin{itemize}
\raggedright
  \item \emph{Say it:} \emph{shimaru}
  \item \emph{Say it:} \emph{shimaru}, once more
  \item \emph{Say it:} \emph{yūbinkyoku wa goji ni shimarimasu} — The post office closes at five.
  \item \emph{Recall:} say the Japanese for to enter, then the Japanese for to be open, then say \emph{shimaru} again
\end{itemize}

\begin{tcolorbox}[breakable,colback=teal!4,colframe=teal!35!black,title={Before you move on}]
What does \ja{しまる} mean? (\textquotedblleft{}To close; to be shut\textquotedblright{}.) Say it once more.
\end{tcolorbox}
\section[\texorpdfstring{\ja{いちば} (ichiba)}{いちば (ichiba)}]{\ja{いちば} (ichiba) — market; marketplace}

\label{lesson:JA-C191-ichiba}



\begin{quote}
Before the new one: say the Japanese for to be open, then the Japanese for to close.
\end{quote}

\subsection*{You'll want to know: \ja{いちば}}
\textbf{\ja{いちば}} — \emph{ichiba} — \textquotedblleft{}market; marketplace\textquotedblright{}.

\begin{quote}
\textbf{\ja{いちばで} \ja{さかなを} \ja{かいます。}} — \emph{ichiba de sakana o kaimasu} — \textquotedblleft{}I buy fish at the market.\textquotedblright{}
\end{quote}

\begin{cousinweb}
\textbf{\ja{いちば}} (\emph{ichiba}) is native Japanese, built from \emph{ichi}, \textquotedblleft{}market\textquotedblright{}, and \emph{ba}, \textquotedblleft{}place\textquotedblright{}, both inherited rather than borrowed (JMdict, KANJIDIC2).
\end{cousinweb}

\begin{grammarlens}[title={What you've built}]
One more word for everyday life at A2.
\end{grammarlens}

\subsection*{Your turn}
\begin{itemize}
\raggedright
  \item \emph{Say it:} \emph{ichiba}
  \item \emph{Say it:} \emph{ichiba}, once more
  \item \emph{Say it:} \emph{ichiba de sakana o kaimasu} — I buy fish at the market.
  \item \emph{Recall:} say the Japanese for to be open, then the Japanese for to close, then say \emph{ichiba} again
\end{itemize}

\begin{tcolorbox}[breakable,colback=teal!4,colframe=teal!35!black,title={Before you move on}]
What does \ja{いちば} mean? (\textquotedblleft{}Market; marketplace\textquotedblright{}.) Say it once more.
\end{tcolorbox}
\section[\texorpdfstring{\ja{さかや} (sakaya)}{さかや (sakaya)}]{\ja{さかや} (sakaya) — liquor store; sake shop}

\label{lesson:JA-C191-sakaya}



\begin{quote}
Before the new one: say the Japanese for to close, then the Japanese for market.
\end{quote}

\subsection*{You'll want to know: \ja{さかや}}
\textbf{\ja{さかや}} — \emph{sakaya} — \textquotedblleft{}liquor store; sake shop\textquotedblright{}.

\begin{quote}
\textbf{\ja{さかやで} \ja{おさけを} \ja{かいます。}} — \emph{sakaya de osake o kaimasu} — \textquotedblleft{}I buy sake at the liquor store.\textquotedblright{}
\end{quote}

\begin{cousinweb}
\textbf{\ja{さかや}} (\emph{sakaya}) is native Japanese, built from \emph{saka-}, the form \emph{sake}, \textquotedblleft{}rice wine\textquotedblright{}, takes in compounds, and \emph{ya}, \textquotedblleft{}house, shop\textquotedblright{} (JMdict, KANJIDIC2).
\end{cousinweb}

\begin{grammarlens}[title={What you've built}]
One more word for everyday life at A2.
\end{grammarlens}

\subsection*{Your turn}
\begin{itemize}
\raggedright
  \item \emph{Say it:} \emph{sakaya}
  \item \emph{Say it:} \emph{sakaya}, once more
  \item \emph{Say it:} \emph{sakaya de osake o kaimasu} — I buy sake at the liquor store.
  \item \emph{Recall:} say the Japanese for to close, then the Japanese for market, then say \emph{sakaya} again
\end{itemize}

\begin{tcolorbox}[breakable,colback=teal!4,colframe=teal!35!black,title={Before you move on}]
What does \ja{さかや} mean? (\textquotedblleft{}Liquor store; sake shop\textquotedblright{}.) Say it once more.
\end{tcolorbox}
\section[\texorpdfstring{\ja{やっきょく} (yakkyoku)}{やっきょく (yakkyoku)}]{\ja{やっきょく} (yakkyoku) — pharmacy; chemist's}

\label{lesson:JA-C191-yakkyoku}



\begin{quote}
Before the new one: say the Japanese for market, then the Japanese for liquor store.
\end{quote}

\subsection*{You'll want to know: \ja{やっきょく}}
\textbf{\ja{やっきょく}} — \emph{yakkyoku} — \textquotedblleft{}pharmacy; chemist's\textquotedblright{}.

\begin{quote}
\textbf{\ja{やっきょくで} \ja{くすりを} \ja{かいます。}} — \emph{yakkyoku de kusuri o kaimasu} — \textquotedblleft{}I buy medicine at the pharmacy.\textquotedblright{}
\end{quote}

\begin{cousinweb}
\textbf{\ja{やっきょく}} (\emph{yakkyoku}) is Sino-Japanese (kango): two Chinese characters read with their on'yomi, readings borrowed from Chinese. They are \emph{yak}, \textquotedblleft{}medicine\textquotedblright{} (Middle Chinese \emph{yak}) and \emph{kyoku}, \textquotedblleft{}office, bureau\textquotedblright{} (\emph{gjowk}) (KANJIDIC2; Middle Chinese in Baxter's transcription).
\end{cousinweb}

\begin{grammarlens}[title={What you've built}]
One more word for everyday life at A2.
\end{grammarlens}

\subsection*{Your turn}
\begin{itemize}
\raggedright
  \item \emph{Say it:} \emph{yakkyoku}
  \item \emph{Say it:} \emph{yakkyoku}, once more
  \item \emph{Say it:} \emph{yakkyoku de kusuri o kaimasu} — I buy medicine at the pharmacy.
  \item \emph{Recall:} say the Japanese for market, then the Japanese for liquor store, then say \emph{yakkyoku} again
\end{itemize}

\begin{tcolorbox}[breakable,colback=teal!4,colframe=teal!35!black,title={Before you move on}]
What does \ja{やっきょく} mean? (\textquotedblleft{}Pharmacy; chemist's\textquotedblright{}.) Say it once more.
\end{tcolorbox}
\section[\texorpdfstring{\ja{みつける} (mitsukeru)}{みつける (mitsukeru)}]{\ja{みつける} (mitsukeru) — to find; to discover}

\label{lesson:JA-C191-mitsukeru}



\begin{quote}
Before the new one: say the Japanese for liquor store, then the Japanese for pharmacy.
\end{quote}

\subsection*{You'll want to know: \ja{みつける}}
\textbf{\ja{みつける}} — \emph{mitsukeru} — \textquotedblleft{}to find; to discover\textquotedblright{}.

\begin{quote}
\textbf{\ja{いちばで} \ja{やすい} \ja{やさいを} \ja{みつけます。}} — \emph{ichiba de yasui yasai o mitsukemasu} — \textquotedblleft{}I find cheap vegetables at the market.\textquotedblright{}
\end{quote}

Polite \textbf{\ja{みつけます}}, \emph{mitsukemasu}; an ichidan verb.

\begin{cousinweb}
\textbf{\ja{みつける}} (\emph{mitsukeru}) is native Japanese, built from \emph{mi}, \textquotedblleft{}seeing\textquotedblright{}, and \emph{tsukeru}, \textquotedblleft{}to attach\textquotedblright{}, both inherited rather than borrowed (JMdict, KANJIDIC2).
\end{cousinweb}

\begin{grammarlens}[title={What you've built}]
One more word for everyday life at A2.
\end{grammarlens}

\subsection*{Your turn}
\begin{itemize}
\raggedright
  \item \emph{Say it:} \emph{mitsukeru}
  \item \emph{Say it:} \emph{mitsukeru}, once more
  \item \emph{Say it:} \emph{ichiba de yasui yasai o mitsukemasu} — I find cheap vegetables at the market.
  \item \emph{Recall:} say the Japanese for liquor store, then the Japanese for pharmacy, then say \emph{mitsukeru} again
\end{itemize}

\begin{tcolorbox}[breakable,colback=teal!4,colframe=teal!35!black,title={Before you move on}]
What does \ja{みつける} mean? (\textquotedblleft{}To find; to discover\textquotedblright{}.) Say it once more.
\end{tcolorbox}
